% Portada de J1.8: la jerarquía de tipos numéricos de Julia y la de un tipo propio, Registro.
\documentclass[tikz,border=0pt]{standalone}
\usepackage{resonancia}
\begin{document}
\begin{tikzpicture}
  \clip[rounded corners=6pt] (0,0) rectangle (12,3.6);
  \shade[top color=rVioleta!14, bottom color=white] (0,0) rectangle (12,3.6);
  \tikzset{abs/.style={draw=rAzul, fill=rAzul!8, rounded corners=2pt, font=\ttfamily\footnotesize, minimum height=0.55cm, inner sep=3pt},
           con/.style={draw=rVioleta, fill=rVioleta!12, rounded corners=2pt, font=\ttfamily\footnotesize, minimum height=0.55cm, inner sep=3pt}}
  \node[abs] (num) at (2.2,3.1) {Number};
  \node[abs] (rea) at (2.2,2.3) {Real};
  \node[abs] (flo) at (1.0,1.5) {AbstractFloat};
  \node[abs] (int) at (3.7,1.5) {Integer};
  \node[con] (f64) at (0.3,0.6) {Float64};
  \node[con] (f32) at (1.9,0.6) {Float32};
  \node[con] (i64) at (3.7,0.6) {Int64};
  \foreach \a/\b in {num/rea, rea/flo, rea/int, flo/f64, flo/f32, int/i64} \draw[eje] (\a) -- (\b);
  \node[abs] (sen) at (8.8,3.1) {Senal};
  \node[con] (reg) at (8.8,2.0) {Registro};
  \draw[eje] (sen) -- (reg);
  \node[etiqueta, anchor=west, font=\ttfamily\scriptsize] at (6.9,1.2) {datos::Matrix{Float64}};
  \node[etiqueta, anchor=west, font=\ttfamily\scriptsize] at (6.9,0.85) {fs::Float64};
  \node[etiqueta, anchor=west, font=\ttfamily\scriptsize] at (6.9,0.5) {canales::Vector{String}};
  \node[etiqueta, color=rAzul] at (5.2,3.1) {\small abstracto};
  \node[etiqueta, color=rVioleta] at (5.2,2.75) {\small concreto};
\end{tikzpicture}
\end{document}
